\begin{pastexam}[label={px:nsysu-110-dmalg-02}]{中山資工所資安組 110 學年度 第 6 題}
（計算）

(Algorithm points) {[10\%]} Find $[37]^{-1}$ in $\boldsymbol{Z}_{66}$ by Euclidean algorithm.
\end{pastexam}

\begin{pastsoln}
答案：$[37]^{-1} = [25]$。 驗算：$37 \times 25 = 925 = 14 \times 66 + 1$，所以 $37 \cdot 25 \equiv 1 \pmod{66}$。

要找 $y$ 使 $37y \equiv 1 \pmod{66}$，也就是找整數 $x, y$ 使 $66x + 37y = 1$。 這樣的 $y$ 存在若且唯若
$\gcd(37, 66) = 1$，而輾轉相除（資結書虛擬碼~6.9）同時算出 $\gcd$ 和 $x$、$y$。

\textbf{輾轉相除}：
\begin{align*}
66 &= 1 \cdot 37 + 29\\
37 &= 1 \cdot 29 + 8\\
29 &= 3 \cdot 8 + 5\\
8 &= 1 \cdot 5 + 3\\
5 &= 1 \cdot 3 + 2\\
3 &= 1 \cdot 2 + 1\\
2 &= 2 \cdot 1 + 0
\end{align*}
最後一個非零餘數是 $1$，$\gcd(37, 66) = 1$，反元素存在。

\textbf{往回代}，每一步把最小的餘數用上一個等式換掉：
\begin{align*}
1 &= 3 - 1 \cdot 2\\
&= 3 - (5 - 3) = 2 \cdot 3 - 5\\
&= 2(8 - 5) - 5 = 2 \cdot 8 - 3 \cdot 5\\
&= 2 \cdot 8 - 3(29 - 3 \cdot 8) = 11 \cdot 8 - 3 \cdot 29\\
&= 11(37 - 29) - 3 \cdot 29 = 11 \cdot 37 - 14 \cdot 29\\
&= 11 \cdot 37 - 14(66 - 37) = 25 \cdot 37 - 14 \cdot 66 .
\end{align*}
兩邊模 $66$：$25 \cdot 37 \equiv 1$，所以 $[37]^{-1} = [25]$。

同樣的計算也可以用表格（每一列維持 $r_i = x_i \cdot 66 + y_i \cdot 37$）：
\begin{center}
\begin{tabular}{rrrrr}
\toprule
$i$ & $r_i$ & $q_i$ & $x_i$ & $y_i$ \\
\midrule
$0$ & $66$ & & $1$ & $0$ \\
$1$ & $37$ & $1$ & $0$ & $1$ \\
$2$ & $29$ & $1$ & $1$ & $-1$ \\
$3$ & $8$ & $3$ & $-1$ & $2$ \\
$4$ & $5$ & $1$ & $4$ & $-7$ \\
$5$ & $3$ & $1$ & $-5$ & $9$ \\
$6$ & $2$ & $1$ & $9$ & $-16$ \\
$7$ & $1$ & $2$ & $-14$ & $25$ \\
\bottomrule
\end{tabular}
\end{center}
$r_{i+1} = r_{i-1} - q_i r_i$，$x$、$y$ 也用同一個式子更新。 最後一列 $1 = -14 \cdot 66 + 25 \cdot 37$，$y = 25$。 若算出的
$y$ 是負的，加上 $66$ 換成 $0$ 到 $65$ 之間的代表元素。
\end{pastsoln}
